% !TEX program = xelatex
\documentclass[a4paper,14pt,oneside,openany]{memoir}

%%% Задаем поля, отступы и межстрочный интервал %%%

\usepackage[left=30mm, right=15mm, top=20mm, bottom=20mm]{geometry}
\pagestyle{plain}
\parindent=1.25cm
\usepackage{indentfirst}

%%% Задаем языковые параметры и шрифт %%%

\usepackage[english,russian]{babel}
\babelfont[russian]{rm}{Times New Roman}
\babelfont[russian]{tt}{Courier New}

%%% Задаем стиль заголовков и подзаголовков в тексте %%%

\setsecnumdepth{subsection}
\renewcommand*{\chapterheadstart}{}
\renewcommand*{\printchaptername}{}
\renewcommand*{\chapnumfont}{\normalfont\bfseries}
\renewcommand*{\afterchapternum}{\hspace{1em}}
\renewcommand*{\printchaptertitle}{\normalfont\bfseries\centering\MakeUppercase}
\setbeforesecskip{20pt}
\setaftersecskip{20pt}
\setsecheadstyle{\raggedright\normalfont\bfseries}
\setbeforesubsecskip{20pt}
\setaftersubsecskip{20pt}
\setsubsecheadstyle{\raggedright\normalfont\bfseries}

%%% Исправление нумерации секций %%%

\renewcommand{\thesection}{\arabic{section}}
\renewcommand{\thesubsection}{\arabic{subsection}}

%%% Сквозная нумерация рисунков и таблиц %%%

\counterwithout{figure}{chapter}
\counterwithout{table}{chapter}

%%% Выравнивание и переносы %%%

\tolerance 1414
\hbadness 1414
\emergencystretch 1.5em
\clubpenalty=10000
\widowpenalty=10000

%%% Задаем параметры оформления рисунков и таблиц %%%

\usepackage{graphicx}
\usepackage{caption}
\usepackage{subcaption}
\usepackage{float}
\graphicspath{{photo/}}
\captionsetup[figure]{font=small, width=\textwidth, name=Рисунок, justification=centering}
\captiondelim{ --- }
\usepackage[section]{placeins}

%%% Задаем параметры ссылок и гиперссылок %%% 

\usepackage{hyperref}
\hypersetup{
    colorlinks=true,
    linkcolor=red,
    citecolor=red
}

%%% Настраиваем отображение списков %%%

\usepackage{enumitem}
\renewcommand*{\labelitemi}{\normalfont{--}}
\setlist{noitemsep, leftmargin=*}

\begin{document}

% Титульный лист (без номера страницы)
\thispagestyle{empty}

\begin{center}
    МИНИСТЕРСТВО ЦИФРОВОГО РАЗВИТИЯ, СВЯЗИ И МАССОВЫХ \\ КОММУНИКАЦИЙ \\ РОССИЙСКОЙ ФЕДЕРАЦИИ

    \vspace{20pt}

    Федеральное государственное бюджетное образовательное учреждение \\ высшего образования \\
    "Сибирский государственный университет телекоммуникаций и \\ информатики"

    \vspace{20pt}

    Кафедра телекоммуникационных систем и вычислительных средств \\ (ТС и ВС)
\end{center}

\vfill

\begin{center}
    Практическая работа №9 \\
    по дисциплине \\
    \textit{"Web-технологии"}

    \vspace{20pt}

    по теме: \\
    \uppercase{Автоматизация сбора данных о клиентах Linux с помощью Ansible}
\end{center}

\vfill

\noindent Студент: \\
\textit{Группа ИКС-532 \hfill А.С. Лёшин}

\vspace{20pt}

\noindent Преподаватель: \\
\textit{старший преподаватель \hfill А.В. Андреев}

\vfill

\begin{center}
    Новосибирск 2026 г.
\end{center}

\newpage
\setcounter{page}{2}
\OnehalfSpacing*

\subsection*{Цель работы}

\begin{itemize}
    \item Настроить Ansible для автоматизированного сбора информации о клиентских машинах.
    \item Собрать с клиентов данные: IP-адрес, версию ОС, имя хоста и свободное место на диске.
    \item Сохранить собранные данные в структурированном виде на сервере.
\end{itemize}

\section{Проверка сетевой связности}

Для выполнения работы используются виртуальные машины:
\begin{itemize}
    \item \textbf{Gateway} (192.168.14.1) — сервер Ansible;
    \item \textbf{Desktop} (192.168.14.11) — клиент 1;
    \item \textbf{WordPress} (192.168.14.6) — клиент 2.
\end{itemize}

На \hyperref[fig-ping]{рисунке \ref*{fig-ping}} представлена проверка связи сервера с клиентскими машинами.

\begin{figure}[H]
    \centering
    \includegraphics[width=0.8\textwidth]{photo/photo1.png}
    \caption{Проверка связи с desktop и wordpress (ping)}
    \label{fig-ping}
\end{figure}

\section{Установка Ansible на сервер}

На сервере Gateway выполняется установка Ansible. На \hyperref[fig-ansible-version]{рисунке \ref*{fig-ansible-version}} представлен вывод команды \texttt{ansible --version} после установки.

\begin{figure}[H]
    \centering
    \includegraphics[width=0.6\textwidth]{photo/photo2.png}
    \caption{Проверка версии Ansible}
    \label{fig-ansible-version}
\end{figure}

\section{Подготовка клиентских машин}

На клиентских машинах установлен и запущен SSH-сервер. На \hyperref[fig-ssh-status]{рисунке \ref*{fig-ssh-status}} показан статус службы SSH на WordPress,На \hyperref[fig-ssh-status1]{рисунке \ref*{fig-ssh-status1}} показан статус службы SSH на Desktop.

\begin{figure}[H]
    \centering
    \includegraphics[width=0.8\textwidth]{photo/photo3.png}
    \caption{Проверка статуса SSH на клиентах (WordPress)}
    \label{fig-ssh-status}
\end{figure}

\begin{figure}[H]
    \centering
    \includegraphics[width=0.8\textwidth]{photo/photo4.png}
    \caption{Проверка статуса SSH на клиентах (Desktop)}
    \label{fig-ssh-status1}
\end{figure}

\section{Настройка SSH-доступа}

На сервере Gateway сгенерирована ключевая пара для аутентификации. Используется алгоритм \texttt{ed25519}, так как он обеспечивает более высокую безопасность при меньшей длине ключа по сравнению с RSA, а также обладает устойчивостью к атакам на основе квантовых вычислений. На \hyperref[fig-ssh-keygen]{рисунке \ref*{fig-ssh-keygen}} показан процесс генерации ключа.

\begin{figure}[H]
    \centering
    \includegraphics[width=0.6\textwidth]{photo/photo5.png}
    \caption{Генерация SSH-ключа на сервере Gateway}
    \label{fig-ssh-keygen}
\end{figure}

Далее ключ копируется на клиентские машины с помощью команды \texttt{ssh-copy-id}. На \hyperref[fig-ssh-copy]{рисунке \ref*{fig-ssh-copy}} показан процесс копирования ключа на оба клиента.

\begin{figure}[H]
    \centering
    \includegraphics[width=0.6\textwidth]{photo/photo6.png}
    \caption{Копирование SSH-ключа на клиентские машины}
    \label{fig-ssh-copy}
\end{figure}

\section{Создание структуры Ansible}

На сервере создаётся файл инвентори \texttt{/etc/ansible/hosts}, в котором перечислены клиенты. На \hyperref[fig-ansible-hosts]{рисунке \ref*{fig-ansible-hosts}} представлено содержимое этого файла.

\begin{figure}[H]
    \centering
    \includegraphics[width=0.6\textwidth]{photo/photo7.png}
    \caption{Содержимое файла инвентори /etc/ansible/hosts}
    \label{fig-ansible-hosts}
\end{figure}

Также создаётся конфигурационный файл \texttt{/etc/ansible/ansible.cfg}, в котором указывается путь к инвентори и отключается проверка host-ключей для автоматизации. На \hyperref[fig-ansible-cfg]{рисунке \ref*{fig-ansible-cfg}} показано его содержимое.

\begin{figure}[H]
    \centering
    \includegraphics[width=0.8\textwidth]{photo/photo8.png}
    \caption{Содержимое файла /etc/ansible/ansible.cfg}
    \label{fig-ansible-cfg}
\end{figure}

\section{Проверка работы Ansible}

Корректность настройки проверяется командой \texttt{ansible clients -m ping}. На \hyperref[fig-ansible-ping]{рисунке \ref*{fig-ansible-ping}} показан результат — оба клиента отвечают успешно.

\begin{figure}[H]
    \centering
    \includegraphics[width=0.8\textwidth]{photo/photo9.png}
    \caption{Проверка доступности клиентов (ansible clients -m ping)}
    \label{fig-ansible-ping}
\end{figure}

\section{Создание и выполнение Playbook}

Создан playbook \texttt{/etc/ansible/playbooks/monitoring.yml}, который собирает информацию о клиентах и сохраняет её в файлы на сервере. На \hyperref[fig-playbook]{рисунке \ref*{fig-playbook}} представлено содержимое playbook.

\begin{figure}[H]
    \centering
    \includegraphics[width=0.8\textwidth]{photo/photo10.png}
    \caption{Содержимое playbook monitoring.yml}
    \label{fig-playbook}
\end{figure}

Запуск playbook выполняется командой \texttt{ansible-playbook}. На \hyperref[fig-playbook-run]{рисунке \ref*{fig-playbook-run}} показан процесс выполнения.

\begin{figure}[H]
    \centering
    \includegraphics[width=0.8\textwidth]{photo/photo11.png}
    \caption{Запуск playbook monitoring.yml}
    \label{fig-playbook-run}
\end{figure}

\section{Результаты сбора данных}

После выполнения playbook в директории \texttt{/etc/ansible/monitoring} созданы файлы с информацией о каждом клиенте. На \hyperref[fig-results]{рисунке \ref*{fig-results}} представлено содержимое файлов \texttt{desktop\_info.txt} и \texttt{wordpress\_info.txt}.

\begin{figure}[H]
    \centering
    \includegraphics[width=0.8\textwidth]{photo/photo12.png}
    \caption{Собранные данные о клиентах}
    \label{fig-results}
\end{figure}

\section*{Заключение}
\addcontentsline{toc}{section}{Заключение}

В ходе выполнения работы были решены следующие задачи:
\begin{enumerate}
    \item Настроена сетевая связность между сервером Gateway и клиентами Desktop и WordPress.
    \item Установлен и настроен Ansible на сервере, настроен SSH-доступ к клиентам.
    \item Созданы файлы инвентори и конфигурации Ansible.
    \item Написан playbook для автоматического сбора информации о клиентах (IP, имя хоста, версия ОС, свободное место на диске).
    \item Playbook успешно выполнен, собранные данные сохранены в директории \texttt{/etc/ansible/monitoring/} на сервере.
\end{enumerate}

В результате автоматизирован процесс сбора диагностической информации о клиентских машинах, что может быть использовано для мониторинга и администрирования сети.

\end{document}